\documentclass{article}
\usepackage{amsmath}
\usepackage{comment}
\usepackage{natbib}
\usepackage{tabularx}

\usepackage{graphicx} % Required for inserting images
\usepackage{xcolor}
\usepackage{algorithm}
\usepackage{algorithmic}
\usepackage{booktabs}
\usepackage{caption}
\usepackage{cancel}
\usepackage{hyperref}
\hypersetup{
   colorlinks=true,
    linkcolor=blue,
  filecolor=magenta,      
    urlcolor=cyan,
citecolor=blue, 
}
\title{Inventory Routing}
\date{March 2024}

\begin{document}

\maketitle
% \section{TO DO}
% \begin{itemize}
%     \item Add a selection mechanism to the local search rather than performing a full enumeration of the neighbourhood.
%     \item Add a disturbance mechanism if required.
%     \item Define parameters used in the algorithm and tune with iRace
%     \item In parallel, documenting the implementation.
% \end{itemize}

\section{Abstract}

The paper presents a multi-start iterated local search  for solving the inventory routing problem with logistic ratio. The proposed optimisation algorithm utilise serval components include an application of the Lagrangian relaxation where violation on the stockout constratins is permitted, a devliery quantity reoptimisation procedure, and serval local search operators integrates both routing decision and quantity decision. The computational experiements were conducted on 978 small instances with 5 to 50 customer and 150 large instances with 50 to 200 customers. The results shown the proposed algorithm outperformance all instances on previous available metaheuristics for large instances and produces competitive results on small instances. We also give our attempt on integrating the hybrid genetic search from capacited vehicle routing problem into the proposed multi-started iterated local search. With hybrid genetic search, the solution quality can still be improved by certain percentage.

\section{Introduction}


The Inventory Routing Problem (IRP) is an extremely challenging optimization problem where the goal is to determine the replenishment and distribution plan for a supply chain network where there is a single supplier and multiple retailers who need to be replenished with the community produced and distributed by the supplier. A planning horizon is defined over which the replenishment activities should take place. Each retailers faces a per period demand over the horizon and has a maximum inventory capacity. The supplier has a fleet of capacitated vehicles which are used for distribution operations. The goal is to define the plan that optimizes a given objective function in such a way that each retailer has a sufficient quantity to satisfy the demand in all periods and the vehicle capacity is not exceeded in any route. More in detail, the plan consists in defining, for each period, which retailers are visited, what is the quantity delivered to each retailer and what are the routes of the vehicles that has to distribute those quantities. 

The objective function which is commonly used in the classical IRP consists of the total cost, to be minimized, which is given by the sum of routing cost plus the inventory cost at the supplier and at the retailers \cite{archetti2012hybrid,archetti2016inventory}. However, it often happens that the inventory cost is hard to define or it is negligible with respect to the routing cost. In this case, other objective functions could be used instead of the total cost. Indeed, a KPI which is often used in practice is the so called \textit{Logistic Ratio} (LR) which consists in the ratio between the routing cost and the total quantity delivered. In this case, we obtain the IRP with LR (IRP-LR). 

The IRP-LR finds relevant applications in practice, especially when the transportation cost is the main component of the total cost. In fact, it was first introduced in the 2016 ROADEF Challenge by Air Liquide. It then gained success and attracted relevant research attention \cite{archetti2017minimizing,archetti2019exact,alvarez2018iterated}. The problem proved to be much more challenging than IRP when it comes to finding effective solution approaches. The goal of the current work is to propose a new heuristic algorithm for the IRP-LR, namely an Iterated Local Search (ILS).

% The inventory routing problem is a well-studied combinatorial optimisation problem for decades. The problem is concerned with a set of customer and a central depot where replenishement schedule is required to avoid customers having stockout respecting the inventory capacity. 

Outline:
\begin{itemize}
    \item As efficient heuristic search procedures exist for routing decisions, while deciding on delivery quantities may be more naturally done using mathematical programming models, it is common in the literature that the IRP is tackled by some matheuristics \cite{iswari2026extending}. 
    \item In \cite{castro2026hybrid}, a hybrid genetic search is designed for the IRP, where it uses hybrid genetic search \cite{vidal2022hybrid} for the routing component, while the delivery quantities are obtained by solving a minimum-cost flow model using the network simplex algorithm.
    \item Allow violation. These two papers allow violation of vehicle capacity 
    \cite{castro2026hybrid}, \cite{archetti2012hybrid}.
\end{itemize}



In the vast literature of IRP, most solution approaches are mathematical formulation-based, mostly because of the delivery quantity decision. A gap remains for a pure heuristic approach.

We assume that a customer will eventually have a stockout. Customer eventually not having a stockout should be discarded. We also assume that there will be no stockout if a customer is visited every day. Meaning that it is not difficult in obtaining a feasible solution. The challenge is to obtain a solution that minimises the transportation cost while delivering as much as possible.

We also assume sufficient production at the depot to fulfil the demand from all retailers.


Recent methodology development focuses on the objective function to minimise the transportation cost and inventory holding cost. \cite{castro2026hybrid,zhao2025large}

\section{Problem Background and Related Literature}
\label{sec:problem_description}

The Inventory Routing Problem (IRP) integrates inventory replenishment and
vehicle-routing decisions over a multi-period planning horizon. Rather than
taking delivery quantities as predetermined inputs, the supplier jointly
determines when customers are visited, how much product is delivered, and how
the corresponding vehicle routes are constructed
\citep{bertazzi2012inventory,coelho2014thirty}. This integration can reduce
system-wide logistics costs, but it also creates strong dependencies between
routing decisions, delivery quantities, and customer inventory levels across
periods.

Most classical IRP formulations minimise an absolute measure of routing,
inventory, or total logistics cost. \cite{archetti2017logistic} introduced an
alternative performance measure, referred to as the \emph{logistic ratio},
defined as the routing cost per unit of product delivered. Their motivation
arises partly from the finite-horizon effect of conventional IRP objectives.
Under total-cost minimisation, there is generally little incentive to deliver
additional inventory near the end of the planning horizon if it is consumed
only after the horizon. A cost-per-unit-delivered objective explicitly
accounts for the amount transported and can therefore mitigate this effect.

For a solution $s$, let $C(s)$ denote its total routing cost and $Q(s)$ the
total quantity delivered. The logistic ratio is

\begin{equation}
    LR(s)=\frac{C(s)}{Q(s)}.
    \label{eq:general_lr}
\end{equation}

If $Q(s)$ were constant across all feasible solutions, minimising
\eqref{eq:general_lr} would be equivalent to minimising routing cost.
In an IRP, however, delivery quantities are endogenous. Consequently, the
objective simultaneously favours

\begin{equation*}
    C(s)\downarrow
    \qquad \text{and} \qquad
    Q(s)\uparrow,
\end{equation*}
and may therefore generate substantially different replenishment decisions
from a conventional cost-minimising IRP.

The problem studied in this paper follows the single-product IRP with logistic
ratio (IRP-LR) introduced in \cite{archetti2017logistic}. The formulation
below additionally introduces auxiliary stockout variables, denoted by
$\mu_{it}$, which are forced to zero in the original problem and therefore do
not modify its feasible solution set. Their purpose is to enable the
Lagrangian-relaxation mechanism used later within the proposed solution
algorithm.


\subsection{Problem Formulation}
\label{sec:problem_formulation}

Let $G=(N,E)$ be a complete undirected graph, where
$N=\{0\}\cup N'$ is the set of locations and $E$ is the set of edges.
Node $0$ represents the depot and
$N'=\{c_1,c_2,\ldots,c_{|N'|}\}$ denotes the set of customers.
A non-negative routing cost $c_{ij}$ is associated with each edge
$\langle i,j\rangle\in E$, and the costs are assumed to satisfy the triangle
inequality. The planning horizon
$T=\{1,\ldots,|T|\}$ consists of discrete periods. Each customer $i\in N'$ has an inventory capacity $U_i$ and consumes
$r_{it}$ units in period $t\in T$. The depot produces $r_{0t}$ units in
period $t$, and $I_{i0}$ denotes the initial inventory at location $i\in N$.
A fleet $K=\{k_1,\ldots,k_{|K|}\}$ of homogeneous vehicles is available,
each with capacity $Q$.

The objective is to construct a distribution plan that minimises the logistic
ratio while ensuring that vehicle capacities and customer inventory
capacities are respected and that no stockout occurs in the original problem.
Split deliveries are not permitted: a customer can be served by at most one
vehicle in each period. The following decision variables are used:

\begin{itemize}
    \item $I_{it}\geq 0$: inventory level at location $i$ at the end of
    period $t$;
    \item $y_{ij}^{kt}$: number of times edge $\langle i,j\rangle$ is
    traversed by vehicle $k$ in period $t$;
    \item $q_{it}^{k}\geq0$: quantity delivered to customer $i$ by vehicle
    $k$ in period $t$;
    \item $z_{it}^{k}\in\{0,1\}$: equal to one if location $i$ is visited
    by vehicle $k$ in period $t$;
    \item $\mu_{it}\geq0$: auxiliary stockout quantity associated with
    customer $i$ in period $t$.
\end{itemize}

For a subset $S\subseteq N'$, define
$E(S)=\{\langle i,j\rangle\in E:i,j\in S\}.
$ 
Assuming that the total delivered quantity is positive, the IRP-LR is
formulated as follows:

\small
\begin{align}
\min \quad &
\frac{
\displaystyle
\sum_{k\in K}
\sum_{\langle i,j\rangle\in E}
\sum_{t\in T}
c_{ij}y_{ij}^{kt}}
{
\displaystyle
\sum_{k\in K}
\sum_{i\in N'}
\sum_{t\in T}
q_{it}^{k}}
\label{obj}
\\
\text{s.t.}\quad
&
I_{0t}
=
I_{0,t-1}+r_{0t}
-
\sum_{k\in K}\sum_{i\in N'}q_{it}^{k},
&& t\in T,
\label{supbal}
\\
&
I_{it}-\mu_{it}
=
I_{i,t-1}-\mu_{i,t-1}-r_{it}
+
\sum_{k\in K}q_{it}^{k},
&& i\in N',\ t\in T,
\label{custbal}
\\
&
I_{it}\geq0,
&& i\in N,\ t\in T,
\label{nonneg_lnv}
\\
&
\mu_{it}\geq0,
&& i\in N',\ t\in T,
\label{nonneg_mu}
\\
&
\sum_{i\in N'}\sum_{t\in T}\mu_{it}\leq0,
&&
\label{zero stockout}
\\
&
\sum_{k\in K}q_{it}^{k}
\leq
U_i-I_{i,t-1},
&& i\in N',\ t\in T,
\label{maxinv}
\\
&
q_{it}^{k}
\leq U_i z_{it}^{k},
&& i\in N',\ k\in K,\ t\in T,
\label{link_qz}
\\
&
\sum_{i\in N'}q_{it}^{k}
\leq Qz_{0t}^{k},
&& k\in K,\ t\in T,
\label{vehcap}
\\
&
\sum_{k\in K}z_{it}^{k}\leq1,
&& i\in N',\ t\in T,
\label{nosplit}
\\
&
\sum_{j:\langle i,j\rangle\in E}y_{ij}^{kt}
=
2z_{it}^{k},
&& i\in N,\ k\in K,\ t\in T,
\label{routing1}
\\
&
\sum_{\langle i,j\rangle\in E(S)}
y_{ij}^{kt}
\leq |S|-1,
&&
S\subseteq N',\ |S|\geq2,\ k\in K,\ t\in T,
\label{subtour}
\\
&
z_{it}^{k}\in\{0,1\},
&& i\in N,\ k\in K,\ t\in T,
\label{binary_z}
\\
&
q_{it}^{k}\geq0,
&& i\in N',\ k\in K,\ t\in T,
\label{nonneg_q}
\\
&
y_{ij}^{kt}\in\{0,1,2\},
&&
\langle i,j\rangle\in E,\ k\in K,\ t\in T.
\label{domain_y}
\end{align}
\normalsize
The objective function~\eqref{obj} minimises total routing cost per unit of
product delivered. Constraints~\eqref{supbal} and~\eqref{custbal} impose the
inventory balance equations at the depot and customers, respectively.
For a customer, $I_{it}-\mu_{it}$ represents its net inventory position.
Since both $I_{it}$ and $\mu_{it}$ are non-negative,
Constraint~\eqref{zero stockout} forces $\mu_{it}=0$ for every customer and
period in the original problem and hence eliminates stockouts.

Constraint~\eqref{maxinv} prevents deliveries from exceeding the available
customer storage capacity, while Constraint~\eqref{link_qz} allows a positive
delivery only if the corresponding customer is visited.
Constraint~\eqref{vehcap} imposes vehicle capacity, and
Constraint~\eqref{nosplit} prevents split deliveries within a period.
Constraints~\eqref{routing1} and~\eqref{subtour} ensure valid vehicle routes.
Because the number of subtour-elimination constraints is exponential,
Constraint~\eqref{subtour} can be generated dynamically when the formulation
is solved exactly. Constraints~\eqref{binary_z}--\eqref{domain_y} specify the
remaining variable domains.


\subsection{Solution Approaches in the Literature}
\label{sec:related_solution_methods}

The fractional objective adds an additional layer of difficulty to the
already challenging interaction between routing, inventory, and replenishment
quantities. The IRP-LR belongs to the broader class of fractional optimisation
problems. The classical parametric approach of \cite{dinkelbach1967}
transforms a fractional problem into a sequence of optimisation problems with
an objective of the form
\[
N(s)-\lambda D(s).
\]
This connection has motivated exact approaches for routing problems with
decision-dependent ratios.

Following the introduction of the IRP-LR,
\cite{archetti2019exact} developed an improved exact algorithm that solves
the fractional problem through a sequence of IRPs with linear objective
functions. \cite{baldacci2020fractional} subsequently investigated
fractional objectives in vehicle routing more generally and developed an
exact framework based on specialised primal and dual bounds.

Metaheuristic approaches are particularly relevant for larger instances.
\cite{alvarez2018iterated} proposed Iterated Local Search (ILS) and Simulated
Annealing algorithms for the IRP and considered the logistic-ratio objective,
providing an important heuristic precedent for the present work. The
ROADEF/EURO Challenge 2016 further stimulated several algorithms for a rich
industrial inventory routing problem with a unit-distribution-cost objective.
These include heuristic branch-cut-and-price
\citep{absi2020branch}, local search combined with mathematical programming
\citep{su2020matheuristic}, fixed-sequence reoptimisation
\citep{he2020fixedsequence}, and hyper-heuristic search
\citep{kheiri2020heuristic}.

Of particular relevance, \cite{he2020fixedsequence} show that substantial
improvements can be obtained by reoptimising delivery quantities and timing
while retaining a fixed customer sequence. This distinction between route
structure and quantity optimisation is especially important under a
logistic-ratio objective because quantities directly affect the denominator
of the objective.

More recently, \cite{alonsopecina2024iterated} proposed an ILS for a
multiple-source inventory-routing application in which operating cost is
divided by total delivered quantity. \cite{man2025logistics} extended
logistics-ratio optimisation to a multi-product setting with lateral
transshipment and uncertain demand.

Recent developments for conventional cost-based IRPs are also methodologically
relevant. \cite{castro2026hybrid} developed a Hybrid Genetic Search (HGS) that
combines population-based search, specialised neighbourhoods, and an
inventory-flow optimisation procedure. \cite{zhao2025large} proposed a large
neighbourhood within HGS that jointly changes visit periods, routing
insertions, and delivery quantities across the planning horizon. Their work
also considers controlled inventory infeasibility through penalised
stockouts. Although these recent HGS approaches do not optimise the same
logistic-ratio formulation considered here, they highlight the potential
benefits of coupling powerful route exploration with specialised quantity
optimisation and of allowing controlled infeasibility during the search.


\begin{table*}[!htbp]
\centering
\caption{Selected literature relevant to the IRP-LR and the proposed
solution methodology.}
\label{tab:literature_irplr}
\scriptsize
\renewcommand{\arraystretch}{1.18}
\setlength{\tabcolsep}{3pt}

\begin{tabularx}{\textwidth}{
>{\raggedright\arraybackslash}p{2.15cm}
>{\raggedright\arraybackslash}p{2.55cm}
>{\raggedright\arraybackslash}p{2.75cm}
>{\raggedright\arraybackslash}X
>{\raggedright\arraybackslash}X}
\toprule
\textbf{Study}
&
\textbf{Problem / objective}
&
\textbf{Solution approach}
&
\textbf{Key search characteristic}
&
\textbf{Relevance}
\\
\midrule

\cite{archetti2017logistic}
&
IRP; routing cost / quantity delivered
&
Exact fractional optimisation
&
Joint routing and replenishment optimisation
&
Introduces and analyses the IRP-LR
\\

\cite{alvarez2018iterated}
&
IRP including logistic-ratio objective
&
ILS and Simulated Annealing
&
Neighbourhood-based routing and replenishment search
&
Direct metaheuristic precedent for IRP-LR
\\

\cite{archetti2019exact}
&
IRP-LR
&
Iterative exact method
&
Sequence of IRPs with modified linear objectives
&
Main exact IRP-LR methodology
\\

\cite{absi2020branch}
&
Rich IRP; unit distribution cost
&
Heuristic branch-cut-and-price
&
Integrated routing, quantity and scheduling optimisation
&
Large-scale fractional industrial IRP
\\

\cite{su2020matheuristic}
&
ROADEF/EURO IRP
&
Local search with MIP/LP
&
Mathematical programming improves timing and quantities
&
Relevant hybrid local-search approach
\\

\cite{he2020fixedsequence}
&
ROADEF/EURO IRP
&
Fixed-sequence matheuristic
&
Reoptimises quantities and timing for fixed visit sequences
&
Closest precedent for route-fixed quantity optimisation
\\

\cite{baldacci2020fractional}
&
Fractional VRP
&
Exact fractional-routing method
&
Primal and dual bounding procedures
&
General fractional-routing methodology
\\

\cite{alonsopecina2024iterated}
&
Operating cost / quantity delivered
&
Iterated Local Search
&
Realistic inventory-routing neighbourhood search
&
Recent ILS with a ratio objective
\\

\cite{man2025logistics}
&
Multi-product logistics-ratio IRP
&
Robust fractional optimisation
&
Uncertainty and lateral transshipment
&
Recent extension of logistics-ratio optimisation
\\

\cite{castro2026hybrid}
&
Conventional IRP
&
Hybrid Genetic Search
&
Route search combined with inventory-flow optimisation
&
Recent high-performance IRP search framework
\\

\cite{zhao2025large}
&
Conventional and stockout-penalised IRP
&
HGS and large-neighbourhood search
&
Cross-period route/quantity optimisation and controlled infeasibility
&
Highly relevant search-design principles
\\

\bottomrule
\end{tabularx}
\end{table*}


\subsection{Lagrangian-Relaxed Problem}
\label{sec:lagrangian_relaxation}

The literature suggests two particularly useful principles for difficult IRPs:
delivery quantities should be reconsidered jointly with the route structure,
and restricting the search to feasible solutions can limit the ability to
move between substantially different replenishment plans. The proposed
algorithm exploits both ideas.

In the original formulation,
Constraint~\eqref{zero stockout}, together with
$\mu_{it}\geq0$, forces all stockout variables to zero. To allow the search
to temporarily explore inventory-infeasible solutions, we dualise
Constraint~\eqref{zero stockout} using a Lagrange multiplier
$\alpha\geq0$. Following the general principle of Lagrangian relaxation
\citep{fisher1981lagrangian}, the resulting relaxed objective is

\begin{align}
\min \quad &
\frac{
\displaystyle
\sum_{k\in K}
\sum_{\langle i,j\rangle\in E}
\sum_{t\in T}
c_{ij}y_{ij}^{kt}}
{
\displaystyle
\sum_{k\in K}
\sum_{i\in N'}
\sum_{t\in T}
q_{it}^{k}}
+
\alpha
\sum_{i\in N'}
\sum_{t\in T}
\mu_{it}.
\label{Relaxed_obj}
\end{align}

The relaxed problem consists of objective~\eqref{Relaxed_obj} subject to

\[
\eqref{supbal}--\eqref{nonneg_mu},
\qquad
\eqref{maxinv}--\eqref{domain_y},
\]

that is, Constraint~\eqref{zero stockout} is removed. The variables
$\mu_{it}$ can therefore take positive values and quantify inventory
shortages, which are penalised through the multiplier $\alpha$.

A small value of $\alpha$ allows the search to explore solutions with larger
stockout violations when this facilitates improvements in routing and
replenishment structure, whereas a larger value increasingly discourages
such violations and directs the search towards the feasible region. Thus,
the relaxed formulation provides a mechanism for controlled exploration of
both feasible and infeasible regions of the solution space.

In what follows, formulation~\eqref{obj}--\eqref{domain_y} is referred to as
the \emph{original problem}, while
\eqref{Relaxed_obj} together with the remaining constraints is referred to as
the \emph{relaxed problem}. The proposed \hyperref[alg:MSILS]{MSILS} algorithm uses the relaxed
problem during local search and dynamically adjusts the multiplier $\alpha$,
as described in Section~\ref{sec: alg}.

\begin{comment}
 \section{Problem description}

Let $G = (N, E)$ be a complete undirected graph where $N = \{0\} \cup N'$ is the set of locations (nodes) and $E$ is the set of edges between locations. Node $0$ represents the depot (supplier location) and $N'=\{c_1,c_2,\dots,c_{|N'|}\}$ denotes the set of customers. A cost $c_{ij}$ is associated with each edge $\langle i, j \rangle$, satisfying the triangle inequality. The planning horizon $T = \{1, \dots, |T|\}$ consists of $|T|$ periods (days). 

Each customer $i \in N'$ has an inventory capacity $U_i$ and a daily demand $r_{it}$ for each period $t \in T$. The quantity produced at the depot during period $t$ is $r_{0t}$, and $I_{i0}$ denotes the initial inventory level at location $i \in N$. A fleet $K=\{k_1,k_2,\dots,k_{|K|}\}$ of homogeneous vehicles, each of capacity $Q$, is available to distribute goods from the depot to the customers.

The Inventory Routing Problem with Logistic Ratio (IRP-LR) aims to determine a distribution plan that minimises the logistic ratio while satisfying all customer demands and the following constraints:

\begin{enumerate}
    \item Vehicle capacity is not exceeded.
    \item No stockout occurs at any node in any period.
    \item The maximum inventory level at customers is not exceeded.
\end{enumerate}

The arc-flow formulation for the IRP-LR requires the following variables:
\begin{itemize}
    \item $I_{it}$: Inventory level at location $i$ at the end of period $t$;
    \item $y_{ij}^{kt}$: Number of times edge $\langle i, j \rangle$ is traversed by vehicle $k$ in period $t$;
    \item $q_{it}^k$: Quantity delivered to customer $i$ by vehicle $k$ in period $t$;
    \item $z_{it}^k$: Binary variable equal to 1 if customer $i$ is visited by vehicle $k$ in period $t$.
    \item \textcolor{red}{$\mu_{it}$: Stockout amount for customer $i$ in period $t$}
\end{itemize}

For a subset of customers $S \subseteq N'$, let $E(S) = \{\langle i, j \rangle \in E \mid i, j \in S\}$. Then, the IRP-LR can be formulated as follows:


\begin{align}
    \min \quad & \frac{\displaystyle \sum_{k \in K} \sum_{\langle i,j\rangle \in E} \sum_{t \in T} c_{ij} y_{ij}^{kt}}
    {\displaystyle \sum_{k \in K} \sum_{i \in N'} \sum_{t \in T} q_{it}^k} \label{obj} \\
    \text{s.t.} \quad 
    & I_{0t} = I_{0,t-1} + r_{0t} - \sum_{k \in K} \sum_{i \in N'} q_{it}^k, && t \in T, \label{supbal} \\
    & I_{it} - \mu_{it} = I_{i,t-1} - \mu_{i,t-1} - r_{it} + \sum_{k \in K} q_{it}^k, && i \in N', \; t \in T, \label{custbal} \\
    & I_{it}  \ge 0, && i \in N, \; t \in T, \label{nonneg_lnv} \\    
    & \mu_{it} \ge 0, && i \in N, \; t \in T, \label{nonneg_mu} \\
    &\sum_{i\in N'}\sum_{t\in T} \mu_{it} \leq 0 \label{zero stockout} \\    
    & \sum_{k \in K} q_{it}^k \le U_i - I_{i,t-1}, && i \in N', \; t \in T, \label{maxinv} \\
    & q_{it}^k \le U_i z_{it}^k, && i \in N', \; k \in K, \; t \in T, \label{link_qz} \\
    & \sum_{i \in N'} q_{it}^k \le Q z_{0t}^k, && k \in K, \; t \in T, \label{vehcap} \\
    & \sum_{k \in K} z_{it}^k \le 1, && i \in N', \; t \in T, \label{nosplit} \\
    & \sum_{j: \langle i,j \rangle \in E} y_{ij}^{kt} = 2 z_{it}^k, && i \in N, \; k \in K, \; t \in T, \label{routing1} \\
    & \sum_{\langle i,j \rangle \in E(S)} y_{ij}^{kt} \le |S| - 1, && S \subset N', \; s \in S, \; k \in K, \; t \in T, \label{subtour} \\
    & z_{it}^k \in \{0,1\}, && i \in N, \; k \in K, \; t \in T, \label{binary_z} \\
    & q_{it}^k \ge 0, && i \in N, \; k \in K, \; t \in T, \label{nonneg_q} \\
    & y_{ij}^{kt} \in \{0,1,2\}, && \langle i,j\rangle \in E, \; k \in K, \; t \in T. \label{domain_y}
\end{align}

The objective function~\eqref{obj} minimizes the \emph{logistic ratio}, defined as the ratio between total routing cost and total quantity delivered.  
Constraints~\eqref{supbal}–\eqref{custbal} represent inventory balance equations for the supplier and customers.  
Constraint~\eqref{nonneg_lnv} ensures that no stockout occurs at any node.  
Constraint~\eqref{maxinv} enforces customer inventory capacity limits.  
Constraint~\eqref{link_qz} links delivery quantities and visit decisions.  
Constraint~\eqref{vehcap} ensures that vehicle capacities are not exceeded.  
Constraint~\eqref{nosplit} prevents split deliveries.  
Constraints~\eqref{routing1}–\eqref{subtour} ensure valid vehicle routing, with subtour elimination enforced dynamically due to its exponential number of constraints.  
Finally, \eqref{binary_z}–\eqref{domain_y} define variable domains. 

By dualising Constraints~\eqref{zero stockout} using Lagrange multiplier $\alpha \geq 0$ gives the following Lagrangian relaxation of the mixed integer linear program \eqref{obj} $-$ \eqref{domain_y}.


\begin{align}
    \min \quad & \frac{\displaystyle \sum_{k \in K} \sum_{\langle i,j\rangle \in E} \sum_{t \in T} c_{ij} y_{ij}^{kt}}
    {\displaystyle \sum_{k \in K} \sum_{i \in N'} \sum_{t \in T} q_{it}^k} + \alpha \sum_{i\in N'}\sum_{t\in T} \mu_{it}  \label{Relaxed_obj} 
\end{align}
\[
\eqref{supbal} -\eqref{nonneg_mu}, \eqref{maxinv}-\eqref{domain_y}
\]

In what follows, this problem will be referred to as the relaxed problem and the problem \eqref{obj}-\eqref{domain_y} will be referred to as the original problem.
\end{comment}

\section{Description of the Algorithm}\label{sec: alg}


This section describes the proposed Multi-Start Iterated Local Search (\hyperref[alg:MSILS]{MSILS}) algorithm for solving the Inventory Routing Problem with Logistic Ratio (IRP-LR). The algorithm uses features that include a mechanism for delivery quantity adjustment, delivery quantity-focused local search operators, and exploring infeasible regions through the amalgamation of iterated local search and Lagrangian relaxation. The remainder of this section is organised as follows. Section \ref{sec:main algorithm} presents the overall \hyperref[alg:MSILS]{MSILS} framework. Section \ref{sec:initial procedure} describes the constructive heuristic used to generate the initial solutions. Section \ref{sec:local search} provides the details on the local search procedure for solving the relaxed problem. Section \ref{sec:rebalance} introduces a quantity-focused rebalancing procedure that attempts to improve total delivery quantity while preserving route schedules. Finally, Section \ref{sec:perturbation} describes the perturbation mechanism used to diversify the search and escape local optima.

\subsection{Main algorithm} \label{sec:main algorithm}


Let $s$ denote a solution, which may be infeasible for the original problem \eqref{obj}-\eqref{domain_y}. For each customer $i\in N'$, each time period $t\in T$, and each vehicle $k\in K$, part of the solution $s$ is characterised by $q^k_{it}(s)$, $I_{it}(s)$, and $\mu_{it}(s)$, where
\begin{itemize}
    \item $q^k_{it}(s)\geq 0$ is the delivery quantity by vehicle $k$, to customer $i$ in time period $t$;
    \item $I_{it}(s) \geq 0$ is the inventory level at customer $i$ in time period $t$; and
    \item $\mu_{it}(s) \geq 0$ is the stockout amount at customer $i$ in time period $t$.
\end{itemize}
If customer $i$ is visited by vehicle $k$ in time $t$, $q^k_{it}(s)$ is positive; otherwise, $q^k_{it}(s)$=0. The net inventory level is defined as $I_{it}(s) - \mu_{it}(s)$ and is computed recursively according to

\begin{equation}\label{initial I0}
I_{i,1}(s) - \mu_{i,1}(s)= I_{i,0} - r_{i,1} + \sum_{k\in K}q^k_{i1}(s),
\end{equation}
\begin{equation}\label{Initial I1}
I_{i,t}(s)-\mu_{i,t}(s)= I_{i,t-1}(s) -\mu_{i,t-1}(s) - r_{i,t} + \sum_{k\in K}q^k_{it}(s),\; 2\leq t\leq |T|.
\end{equation}
If $s$ is feasible to the relaxed problem and has $\sum_{i\in N'} \sum_{i\in T} \mu_{it}(s) = 0$, then $s$ is also feasible for the original problem. The value of $s$ on the objective function \eqref{obj} is denoted by $f(s)$, and the value on the objective function \eqref{Relaxed_obj} with Lagrange multiplier $\alpha$ is denoted by $f_{\mu}(s, \alpha)$. The following pseudocode outlines the proposed \hyperref[alg:MSILS]{MSILS}. 



\begin{algorithm}[H]
\caption{\textbf{MSILS}}\label{alg:MSILS}
\begin{algorithmic}[1]
\STATE $s^* \gets$ Empty \AND  $f(s^*) \gets +\infty$  \label{initialisation}
\FOR{$i\gets 1$ \TO $Max_{restart}$} \label{start msils}
\STATE $s\gets$ \hyperref[alg:InitialSchedule]{Construct initial solution} \label{initial solution construction}
\IF{$f(s)<f(s^*)$} \STATE $s^*\gets s$ \ENDIF
\STATE  $it \gets 0$ \label{initilise eta}
\WHILE{ $it < Max_{it}$} \label{start ils}
    \STATE $\alpha \gets \epsilon (s^*, s , 0, \eta)$ \label{LS:initialise alpha}  
    \STATE $\bar{s} \gets s$
    \STATE $s\gets$ \hyperref[alg:LocalSearch]{Local\_Search}$(s,\alpha)$    
    \WHILE{$f_{\mu}(s,\alpha) \neq f_{\mu}(\bar{s},\alpha)$ \OR $s$ has stockout} \label{search: repeat}
        \STATE  $\alpha \gets \epsilon (s^*, s , \alpha, \eta)$\label{LS:update alpha}        
        \STATE $\bar{s} \gets s$
        \STATE $s\gets$\hyperref[alg:LocalSearch]{Local\_Search}$(s,\alpha)$       
    \ENDWHILE \label{search: until}
    \IF{$s$ is feasible \AND $f(s)<f(s^*)$}
            \STATE $s^*\gets s$, $it\gets 0$
            \STATE $s^*\gets$ \hyperref[alg:rebalance]{Rebalance}$(s^*)$ \label{run rebalance}
    \ENDIF
    \STATE $s \gets$ Disturbance$(s^*, it)$, $it\gets it+1$ \label{disturb}
\ENDWHILE \label{end ils}
\ENDFOR \label{end msils}
\RETURN solution $s^*$
\end{algorithmic}
\end{algorithm}



% \begin{algorithm}[H]
% \caption{\textbf{MSILS}}\label{alg:MSILS}
% \begin{algorithmic}[1]
% \STATE $s^* \gets$ Empty \AND  $f(s^*) \gets +\infty$  \label{initialisation}
% \FOR{$i\gets 1$ \TO $Max_{restart}$} \label{start msils}
% \STATE $s\gets$ \hyperref[alg:InitialSchedule]{Construct initial solution} \label{initial solution construction}
% \IF{$f(s)<f(s^*)$} \STATE $s^*\gets s$ \ENDIF
% \STATE  $it \gets 0$, $\eta \gets \eta'$, $u\gets 0$ \label{initilise eta}
% \WHILE{ $it < Max_{it}$} \label{start ils}
%     \STATE $\alpha \gets \epsilon (s^*, s , 0, \eta)$ \label{LS:initialise alpha}  
%     \STATE $\bar{s} \gets s$
%     \STATE $s\gets$ \hyperref[alg:LocalSearch]{Local\_Search}$(s,\alpha)$    
%     \WHILE{$f_{\mu}(s,\alpha) \neq f_{\mu}(\bar{s},\alpha)$ \OR $s$ has stockout} \label{search: repeat}
%         \STATE  $\alpha \gets \epsilon (s^*, s , \alpha, \eta)$\label{LS:update alpha}        
%         \STATE $\bar{s} \gets s$
%         \STATE $s\gets$\hyperref[alg:LocalSearch]{Local\_Search}$(s,\alpha)$       
%     \ENDWHILE \label{search: until}
%     \IF{$min\{it, u\} \geq U$}
%         \STATE $u\gets 0$, $\eta \gets \kappa \times \eta$ \label{update eta}
%     \ENDIF
%     \IF{$s$ is feasible \AND $f(s)<f(s^*)$}
%             \STATE $s^*\gets s$, $it\gets 0$
%             \STATE $s^*\gets$ \hyperref[alg:rebalance]{Rebalance}$(s^*)$ \label{run rebalance}
%     \ENDIF
%     \STATE $s \gets$ Disturbance$(s^*, it)$, $it\gets it+1$, $u\gets u+1$ \label{disturb}
% \ENDWHILE \label{end ils}
% \ENDFOR \label{end msils}
% \RETURN solution $s^*$
% \end{algorithmic}
% \end{algorithm}

Initially, the incumbent best solution $s^*$ is empty, and $f(s^*)$ is set to $+\infty$ (line \ref{initialisation}). Throughout the search process, whenever a feasible solution to the original problem is found with an objective value lower than $f(s^*)$, the incumbent solution $s^*$ is updated accordingly. The FOR loop (lines \ref{start msils}-\ref{end msils}) implements the multi-start framework. At each restart, a feasible solution to the original problem is generated using the constructive heuristic described in \ref{sec:initial procedure}. This solution is used as the starting solution for the upcoming application of the iterated local search (WHILE loop, lines \ref{start ils}-\ref{end ils}). By repeatedly restarting the iterated local search from different initial solutions, the algorithm can explore diversified regions of the solution space. The number of starting solutions is controlled by the parameter $Max_{restart}$.

Each iteration of the iterated local search uses a local search procedure to solve a sequence of relaxed problems, each with a different value of $\alpha$. If an optimal solution $\hat{s}$ for the relaxed problem can be obtained, then following a commonly used version of the Lagrangian relaxation method \cite{fisher1981lagrangian}, $\alpha$ is computed by 

\begin{equation}
\alpha + \frac{\eta \times (f(s^*) - f_{\mu}(\hat{s},\alpha)) \times (\sum_{i\in N'}\sum_{t\in T} \mu_{it}(\hat{s}) )\,} {\left(\sum_{i\in N'}\sum_{t\in T} \mu_{it}(\hat{s}) \right)^2}  \label{classic lr formula}
\end{equation}
where $\eta$ is a scalar. Since the local search cannot guarantee the optimal solution $\hat{s}$, instead of \eqref{classic lr formula}, $\alpha$ is computed by 
\begin{equation}
\epsilon(s^*, s, \alpha, \eta) = \alpha + \frac{\eta \times \delta \times (\sum_{i\in N'}\sum_{t\in T} \mu_{it}(s) + 0.001)\,} {\left(\sum_{i\in N'}\sum_{t\in T} \mu_{it}(s) \right)^2+ 0.00001}  \label{epsilon}
\end{equation} where
\begin{equation*}    
\delta= \left\{ \begin{array}{ll}
  		f(s^*) - f_{\mu}(s,\alpha)  & \mbox{if $f(s^*) - f_{\mu}(s,\alpha) >0$}; \\
 		 0.1\times f(s^*)& \mbox{ Otherwise.} 
  \end{array}\right. 
\end{equation*}
Once local search obtains a feasible solution for the original problem, $\epsilon(\cdot)$ becomes just 
\begin{equation*}
\alpha + \frac{\eta \times \delta \times ( 0.001)\,} { 0.00001} 
\end{equation*} 
 which greedily increases $\alpha$ until the stopping criterion is satisfied for the WHILE loop (lines \ref{search: repeat}-\ref{search: until}). The initial value of $\alpha$ is computed by $\epsilon (s^*, s, 0, \eta)$ in line \ref{LS:initialise alpha} and the subsequent values is computed by $\epsilon (s^*, s, \alpha, \eta)$ in line \ref{LS:update alpha} using $\alpha$ from the previous call of the local search procedure.  

To escape from the local optimum, the \hyperref[alg:MSILS]{MSILS} begins with a quantity-focused disturbance procedure. This procedure referred to as the rebalance procedure attempts to adjust the delivery quantities on the incumbent best solution, keeping the schedule fixed (see Section \ref{sec:rebalance}). To limit computational overhead, the rebalance procedure is invoked only when the incumbent best solution has been improved during the current local search iteration. To further diversify the search, the \hyperref[alg:MSILS]{MSILS} applies a route-focused disturbance procedure on the incumbent best solution, generating the starting solution for the subsequent call of local search. The iterated local search (WHILE loop lines \ref{start ils}-\ref{end ils}) terminates after the counter $it$ reaches the maximum number $Max_{it}$. Various parameter values for $Max_{restart}$, $Max_{it}$, and $\eta$ are examined in Section \ref{sec: computational experiements}.



\subsection{Initial feasible solution} \label{sec:initial procedure}

This section describes the constructive heuristic used to generate a feasible solution to the problem \eqref{obj}-\eqref{domain_y}. The aim of this procedure is to generate good starting solutions for the \hyperref[alg:MSILS]{MSILS} and to establish the initial upper bound. The proposed heuristic incorporates the routing-related considerations into the construction of replenishment schedules by determining customers' priority based on both replenishment urgency and their distance from the depot.

Let $s$ denote a solution. It is initially empty, with an empty schedule where routes are empty, and the delivery quantities are set to zero, i.e., $q^k_{it}(s)=0$ for $i\in N'$, $k\in K$, and $t\in T$. The net inventory levels $I_{it}(s) - \mu_{it}(s)$ are then calculated recursively by \eqref{initial I0} and \eqref{Initial I1}. Then, the customers are partitioned by the earliest time period $t$ at which $I_{it}(s) - \mu_{it}(s)$ becomes negative (i.e., when a stockout occurs). Such a time reflects the relationship between a customer's demand and its inventory capacity. Customers with lower inventory capacity or higher demand will experience stockout earlier and consequently require more frequent replenishments. Therefore, customers with earlier stockout times are assigned a higher priority.

Furthermore, within each priority group corresponding to time period $t$, customers are further partitioned spatially based on their locations. For each customer $i$ in $N'$, its location is specified by coordinates $(x_i, y_i)$. Then, let 
$$x_{min} = \min\{x_i:i\in N'\}, x_{max} = \max\{x_i:i\in N'\},$$ and
$$y_{min}= \min\{y_i:i\in N'\}, y_{max}= \max\{y_i:i\in N'\}.$$
All customers lie within the axis-aligned rectangular region determined by the corner points $(x_{min},y_{min})$, $(x_{max},y_{min})$, $(x_{max},y_{max})$, $(x_{min},y_{max})$. This region is partitioned into a grid of $d\times d$ equally sized subregions, denoted by $R=\{r_1,r_2,\dots,r_{d\times d}\}$. To enable diversified solutions, the value of $d$ is randomly selected from $\left\{ d_{\min} + 2m: m = 0,\ldots,d_{\max} \right\}$ where $d_{min}$ and $d_{\max}$ are parameters defining the minimum and maximum possible grid resolution. The subregions in $R$ are ordered in descending order of their distance from the depot, where the distance is measured between the depot and the centre of each subregion. 

This ordering is motivated by the observation that, during the early stages of the construction, vehicles typically have sufficient residual capacity to deliver large quantities while respecting both vehicle and inventory capacity constraints. Such deliveries reduce the frequency with which distant customers need to be visited throughout the planning horizon. In contrast, customers considered later in the construction process may receive smaller deliveries because vehicle capacity has already been allocated to previously scheduled customers, potentially resulting in more frequent visits. Since repeated visits to distant customers can substantially increase transportation costs, customers located farther from the depot are prioritised. No further routing optimisation is performed during the constructive phase; instead, routing improvements are devoted to the local search procedure described in Section \ref{sec:local search}.

Let $P_{\rho,r} \subseteq N'$ denote the subset of customers in priority group $\rho\in T$ and subregion $r\in R$. The pseudocode, \hyperref[alg:InitialSchedule]{Construct Initial Solution}, summarises the procedure.


\begin{algorithm}[H]
\captionsetup{labelformat=empty}
\caption{\textbf{Construct Initial Solution}}
\label{alg:InitialSchedule}
\begin{algorithmic}[1]
\STATE initialise solution $s$ \label{Initialise s} 
\STATE partition customers by priority and distance to depot
\FOR{$\rho \in T$ \AND $r\in R$}
    \WHILE{$P_{\rho,r}\neq [\hspace{0.1cm}]$}
        \STATE randomly choose a customer $i$ from $P_{\rho,r}$
        \STATE $t_{current}\gets 1$
        \WHILE{ $\{t\in T: I_{i,t}(s)-\mu_{i,t}(s) <0\} \neq [\hspace{0.1cm}]$}
            \STATE $t_{stockout}\gets \min\{t\in T: I_{i,t}(s)-\mu_{i,t}(s) <0\}$
            \STATE $t_{assign}\gets$ \hyperref[alg:FindVisitTime]{Find Feasible Visit Time} ($t_{current}, t_{stockout}, i$)
            \STATE $(k,q)\gets $ \hyperref[alg:FindVehicleToDeliverMax]{Find Vehicle To Deliver Max}($t_{assign}, i, s$)            
            \STATE $q^k_{i,t_{assign}}\gets q$
            \STATE add customer $i$ to vehicle $k$'s route on time period $t_{assign}$
            \STATE for all $t'\in T$, update $I_{i,t'}(s)- \mu_{it'}(s) $ according to \eqref{initial I0} and \eqref{Initial I1} 
            \STATE $t_{current}\gets t_{assign}+1$
        \ENDWHILE
        \STATE $P_{t,r}\gets P_{t,r}\setminus i$
    \ENDWHILE
\ENDFOR
\STATE $s\gets $
\RETURN $s$ 
\end{algorithmic}
\end{algorithm}

The heuristic constructs a feasible solution by iteratively inserting customer visits and delivery quantities until there are no more stockouts. More specifically, for each subset $P_{\rho,r}$, the heuristic randomly selects a customer $i$ and repeatedly resolves its earliest remaining stockout. At each iteration, $t_{stockout}$ is set to the first period in which $I_{i,t}(s)-\mu_{i,t}(s)<0$. Then, the heuristic determines a suitable visit period $t_{assign}$ by \hyperref[alg:FindVisitTime]{Find Feasible Visit Time} and determines a suitable visit vehicle $k$ and delivery quantity $q$ by \hyperref[alg:FindVehicleToDeliverMax]{Find Vehicle To Deliver Max}. The resulting net inventory levels are updated according to \eqref{initial I0} and \eqref{Initial I1}. The process continues until no stockout remains for this customer. Then, this customer is removed from the current subset. After customers in every subset $P_{\rho,r}$ have been processed, the output solution $s$ is returned to \hyperref[alg:MSILS]{MSILS}. If a customer still experiences a stockout even when visited in every period of the planning horizon, then no feasible schedule exists for that customer. In this case, the instance is infeasible. The procedure returns an empty solution, and the \hyperref[alg:MSILS]{MSILS} is terminated. 




% Minimum total delivery quantity for customer $i\in N'$ $= \sum^{T}_{t=1}d_{i,t}-I_{i,0}$

% Approximate number of visits $v_i= \frac{\sum^T_{t=1}d_{i,t}-I_{i,0}}{U_i}$, $v_i= \frac{\sum^T_{t=1}d_{i,t}-I_{i,0}}{\min{I_{i,t}|}}$

% Minimum delivery quantity per visit $=\frac{\sum^T_{t=1}d_{i,t}-I_{i,0}}{v_i}$








\subsection{Local search} \label{sec:local search}

% From the vast literature of vehicle routing and inventory routing, there have been validated that permitting some constraint violation in local search

The local search uses five neighbourhood operators, denoted by $O_1$, $O_2$, $O_3$, $O_4$, and $O_5$.  These operators are based on those commonly found in the field of vehicle routing (see, for example, \cite{kindervater1997}, \cite{laporte2000classical}, and \cite{elshaer2020taxonomic}) and are also frequently used in the field of inventory routing (see, for example, \cite{alvarez2018iterated}, \cite{lagana2024iterated}, and \cite{castro2026hybrid}). Each operator defines a neighbourhood through a set of alterations. Each alteration transforms the input solution into a neighbouring solution by modifying one or more of its components.  

Rather than employing separate operators for route alterations and delivery quantity alterations, the operators $O_1$, $O_2$, $O_3$, and $O_4$ integrate both into a single alteration. Such an integration enables the exploration of the logistic ratio, i.e., the simultaneous exploration of total transportation cost and total delivery quantity. The routing component of these operators is described as follows:


\begin{itemize}
\item Operator $O_1$ moves up to 3 adjacent customers from a route to another in the same period. 
\item Operator $O_2$ exchanges up to 2 adjacent customers in a route with up to 2 adjacent customers in another route in the same period.
\item Operator $O_3$ transfers a customer visit to a different time period in which this customer is not visited, with the vehicle giving the largest possible delivery quantity at a position with the lowest in transportation cost.
\item Operator $O_4$ exchanges a customer in a route with one unvisited customer, including the case where either a visited customer becomes unvisited or an unvisited customer becomes visited.
\end{itemize}

% \item Things worth mentioning: \textcolor{red}{In \cite{alvarez2018iterated}, $Op-opt, Shift, Swap$ are routing operators. So it is reasonable to say that there is no update for the delivery quantities and inventory levels for the subsequent days. Also, I suspect any infeasible move to an infeasible solution is skipped. In our operators, for $Or-opt(k), k \in \{1, 2\}$, it is pure routing operators. But for $Shift(k), k \in \{1, 2\}$ and $Swap(k1, k2), k1, k2 \in \{1, 2\}, k1 \geq k2$, after the move, delivery quantities and inventory level are updated.}


For an alteration that results in removing the visit to customer $i\in N'$ at time $t\in T$ on vehicle $k\in K$, the associated $q^k_{it}(s)$ becomes 0. For an alteration that results in a new visit to customer $i\in N'$ at time $t\in T$ on vehicle $k\in K$, the associated $q^k_{it}(s)$ is computed by
\begin{equation}
   q^k_{jt} = \max\{0,\min\{Q - \sum_{i\in N'}q^k_{it}(s), U_i-I \}\}\label{Adjust quantity}
\end{equation}
where 
\begin{equation*}    
I= \left\{ \begin{array}{ll}
  		I_{i,0}  & \mbox{if $t =1$}; \\
 		I_{i,t-1}(s) & \mbox{ if $2\leq t\leq |T| $} 
  		\end{array}\right. 
\end{equation*}
Please note that $\max\{0,\cdot\}$ is required since the local search allows the violation of stockout where $I$ can be negative and $U_i-I$ can also be negative. For customer $i$, the resulting inventory level is updated recursively according to \eqref{initial I0} and \eqref{Initial I1}. These updated inventory levels may affect the subsequent deliveries. Therefore, following the update in each time period, if the customer is visited in the subsequent time period, the delivery quantity for that visit is adjusted using \eqref{Adjust quantity}.

For the last operator, Operator $O_5$ is a pure routing operator. It transfers up to 3 adjacent customers from their current position to another position in the same route. In contrast to operators $O_1$, $O_2$, $O_3$, and $O_4$, the alterations associated with operator $O_5$ do not affect the delivery quantities and the inventory levels.

\begin{algorithm}
\captionsetup{labelformat=empty}
\caption{\textbf{Local Search$(s,\alpha)$}} \label{alg:LocalSearch}
\begin{algorithmic}[1]
\FOR{$i \gets 0;\ i \leq 5;\ i \gets i+1$} \label{search: for}
    \REPEAT \label{search: repeat N}
        \STATE $\hat{s} \gets s$, $j \gets 1$
        \WHILE{$j < |O_i(s)|$} 
            \STATE $s' \gets s_j \in O^i(s)$,  $j \gets j+1$
                \IF{$f_{\mu}(s',\alpha) < f_{\mu}(s,\alpha)$}
                     \STATE $s \gets s'$
                \ENDIF
        \ENDWHILE
    \UNTIL{$f_{\mu}(s,\alpha) = f_{\mu}(\hat{s}, \alpha)$} \label{search: until N}        
\ENDFOR \label{search: endfor}
\STATE $s\gets $ remove any visit with zero delivery quantity
\RETURN $s$
\end{algorithmic}
\end{algorithm}


The local search procedure implements a sequential variable neighbourhood descent procedure for solving the relaxed problem. For a given solution $s$ and the Lagrange multiplier $\alpha$, the neighbourhoods induced by operators $O_1$ to $O_5$ are explored successively. Within a neighbourhood associated with Operator $i\in \{1,2,3,4,5\}$, all neighbouring solutions denoted by $O_i(s)$ are evaluated according to the objective function $f_{\mu}(\cdot,\alpha)$, and the current solution is updated whenever an improving neighbour is identified. The neighbourhood is repeatedly re-explored until no improving solution can be found, thereby ensuring local optimality with respect to the corresponding operator. The search then proceeds to the next operator. The procedure terminates once all five operators have been exhausted without improvement, at which point any visits associated with zero delivery quantity are removed and the resulting solution is returned. The pseudocode, \hyperref[alg:LocalSearch]{Local Search} summaries the local search procedure.



\subsection{Rebalance Procedure} \label{sec:rebalance}

The local search procedure simultaneously explores routing decisions and delivery quantities. However, the improved solution may still contain delivery quantities that are not fully consistent with the fixed schedule. In particular, some customer visits may carry quantities that are unnecessarily large, whereas other visits may have unused inventory or vehicle capacity that could be exploited to increase the total delivery quantity. Since the objective of the IRP-LR is expressed as the ratio between routing cost and delivery quantity, improving the quantity allocation alone can further improve the objective value even when the routing schedule remains unchanged. This section describes the rebalance procedure, which exploits such an opportunity whenever a new incumbent solution $s^*$ is obtained. 

The procedure consists of two phases. First, delivery quantities are partially removed from existing visits whenever such reductions do not create stockouts. This releases inventory and vehicle capacity that may have been inefficiently allocated. Second, the released quantities are reallocated among the existing visits. Customers are selected randomly, favouring customers with low inventory levels, high demand rates, long intervals until their next visits, and infrequent replenishment patterns. For each selected customer, the maximum feasible quantity is assigned while respecting inventory capacity and vehicle capacity. The inventory level is then updated, and visits associated with zero delivery quantity are subsequently removed.

By combining quantity reallocation with randomised customer selection, the rebalance procedure not only intensifies the search for better delivery quantities around a promising routing schedule but also introduces diversification, enabling the algorithm to escape local optima in the quantity domain. The pseudocode below summarises the rebalance procedure.

\begin{algorithm}[H]
\caption{\textbf{Rebalance$(s^*)$}}\label{alg:rebalance}
\begin{algorithmic}[1]
\STATE $s \gets s^*$
\FOR{$i \gets 1$ \TO $|N'|$}     
    \FOR{$t \gets 1$ \TO $|T|$} 
        \IF{customer $i$ is visited on day $t$}
            \STATE Quantity $\gets $ Quantity - \hyperref[alg:FeasibleQuantityRemoval]{Feasible\_Quantity\_To\_Removal}$(i, t)$
            \STATE Update inventory levels for all subsequent days
        \ENDIF
    \ENDFOR
\ENDFOR
\STATE $s\gets$\hyperref[alg:AssignQuantity]{Assign\_max\_possible\_quantity\_after\_feasible}$(s)$
\STATE $s\gets$Remove visits with zero delivery quantities
\RETURN $s$
\end{algorithmic}
\end{algorithm}




In addition to the rebalance procedure described above, an alternative \hyperref[alg: destroy and rebalance]{Destory\_And\_Rebalance} procedure was developed. This variant first removes all delivery quantities and then reconstructs them while keeping the routing schedule fixed. Because the reconstruction does not guarantee feasibility, a repair mechanism is subsequently applied to eliminate stockouts whenever possible. The larger perturbation induced by this procedure promotes diversification and allows the search to explore quantity allocations that are less accessible through the former rebalance procedure. The pseudocode for this variant is presented in Appendix~\ref{sec: rebalance without feasibility}. Both rebalance procedures are examined in the computation experiments.


\subsection{Perturbation} \label{sec:perturbation}


The perturbation procedure is a diversification mechanism designed to help the \hyperref[alg:MSILS]{MSILS} escape local optima. It randomly removes two or more adjacent customers from a route together with their associated delivery quantities, after which inventory levels and subsequent deliveries are updated accordingly. The number of removals is determined by $1+\lfloor\frac{|K|\times it}{Max_{it}}\rfloor$, where $it$ is the number of consecutive non-improving iterations, $Max_{it}$ is the maximum number of consecutive non-improving iterations allowed, and $|K|$ is the number of vehicles available in a period. Hence, the perturbation strength increases progressively. This adaptive strategy allows the algorithm to balance intensification and diversification by exploring increasingly distant regions of the solution space when improvements become difficult to obtain. The resulting perturbed solution serves as the starting point for the next local search iteration. Since the removals typically induce stockouts, the perturbed solution is generally infeasible for the original problem. However, allowing such infeasibility enables the search to traverse infeasible regions of the solution space and escape deep local optima. Feasibility is subsequently restored during the local search.


\section{Computational experiments} \label{sec: computational experiements}

This section describes the computational experiments conducted to evaluate the performance of the proposed \hyperref[alg:MSILS]{MSILS}. The algorithm is implemented in C++, and all experiments were conducted on a computer with an Intel Core i7-1185G7 3.00 GHz CPU and 16GB RAM. The rest of the section is structured as follows. The test instances used for the computation experiments are summarised in Section \ref{sec: test instances}. Section \ref{sec: tuning} discusses the parameter settings and demonstrates the performance gain from different components used in the \hyperref[alg:MSILS]{MSILS}. Section \ref{sec:performance evaluation} evaluates the \hyperref[alg:MSILS]{MSILS} by comparing its performance with the performance of the ILS and SA algorithms proposed by \cite{alvarez2018iterated}. This section also evaluates the performance gain from integrating the Hybrid Genetic Search (HGS) framework of \cite{vidal2022hybrid} into \hyperref[alg:MSILS]{MSILS}.

\subsection{Test instances} \label{sec: test instances}

To evaluate the proposed \hyperref[alg:MSILS]{MSILS}, computational experiments were conducted on benchmark instances originating from \cite{archetti2007branch}, \cite{archetti2012hybrid}, and the 12th DIMACS Implementation Challenge for the multi-vehicle IRP. These instances have been widely used in the IRP literature, including \cite{archetti2012hybrid, coelho2013branch, desaulniers2016branch, santos2016hybrid, archetti2017logistic, alvarez2018iterated,  archetti2019exact, castro2026hybrid}.

The benchmark introduced by \cite{archetti2007branch} contains a set of single-vehicle IRP instances comprising 50 instances ranging from 5 to 50 customers for a planning horizon of 3 time periods, and 30 instances ranging from 5 to 30 customers for a planning horizon of 6 time periods. Each instance is available under two inventory-cost scenarios (high and low), resulting in a total of 160 instances. These single-vehicle IRP instances were extended to multi-vehicle versions with a number of vehicles ranging from 2 to 5, creating an additional 640 instances. Following \cite{alvarez2018iterated} and the DIMACS Implementation Challenge, the vehicle capacities for the multi-vehicle IRP instances were obtained by dividing the vehicle capacity in the corresponding single-vehicle instances by the number of vehicles and rounding down to the nearest lower integer.

The 12th DIMACS Implementation Challenge further introduced 160 additional multi-vehicle instances with a number of customers ranging from 35 to 50, completing the 6-time-period planning horizon dataset. In addition, the single- and multi-vehicle IRP instances for 3 planning horizons were extended to a planning horizon of 4 and 5 periods in \cite{archetti2017logistic}. This was done by increasing the planning horizons from 3 to 4 and 5, resulting in another 1000 instances. Collectively, these datasets comprise 1960 instances covering 5 to 50 customers, 3 to 6 planning horizons, 1 to 5 vehicles, and two inventory-cost settings. 

A larger scale of the benchmark was also considered. In \cite{archetti2012hybrid}, a set of 60 single-vehicle IRP instances was introduced with a number of customers ranging from 50 to 200. These instances were also extended to multi-vehicle instances, resulting in a total of 300 instances. The single-vehicle IRP instances from \cite{archetti2007branch} and \cite{archetti2012hybrid} can be downloaded from \href{http://www.leandro-coelho.com/instances/inventory-routing/}{Single-vehicle IRP instances} and the instances for the 12th DIMACS Implementation Challenge for the multi-vehicle IRP can be downloaded from \href{https://dimacs.rutgers.edu/programs/implementation-challenges/vehicle-routing/implementation-challenge-vehicle-routing/inventory-routing}{DIMACS Inventory Routing Challenge}. 

For the above-mentioned instances, the distance $c_{ij}$ between each node $i\in N$ and $j\in N$ is calculated and rounded to the nearest integer by \[ \left\lfloor \sqrt{(x_i-x_j)^2 + (y_i-y_j)^2} +\frac{1}{2} \right\rfloor \], where $x_i,x_j,y_i,y_j$ are the x- and y-coordinates for nodes $i$ and $j$. Since the IRP-LR does not include inventory cost, only one inventory-cost setting was retained, and inventory-cost information was discarded. Consequently, the final test bed consists of 980 instances for 5 to 50 customers and 150 instances for 50 to 200 customers. 

Also noting, two instances were excluded from the study due to infeasibility. The instance ``abs5n5\_5" with 6 periods and 5 vehicles is a known infeasible instance \cite{alvarez2018iterated}. In addition, the preliminary experiments revealed that the instance ``abs4n5\_5" with 5 planning periods and 5 vehicles is also infeasible, as stockouts occur even when a certain customer is visited in every period.

For reproducibility, all benchmark instances used in the computational experiments are publicly 
available at \url{https://github.com/AndyElhamSarang/IRP_instances.git}.


\subsection{Impact on the performance of individual mechanisms in \hyperref[alg:MSILS]{MSILS} and parameter tuning} \label{sec: tuning}



The parameters for the \hyperref[alg:MSILS]{MSILS} are set to the values listed in Table \ref{default parameters}. These values were determined from preliminary experiments.

\begin{table}[H]
\caption{Default values of parameters}
\label{default parameters}
\centering
\begin{tabular}{ p{2cm}p{5cm}p{7cm}}
\toprule%
\textbf{Parameters} & \textbf{Recommended Value} &\textbf{Tested Range} \tabularnewline
\hline
$d_{\min}$&2& $\{1,2,3,4,5\}$ \tabularnewline
$d_{\max}$&1& $\{0,1,2,5\}$ \tabularnewline
$Max_{restart}$ & 3 & \{1,3,5,10,20\} \tabularnewline
$U$ & $5$ & \{ 1,5,10,30,50\}\tabularnewline
$\kappa$ & $1.0$ & \{ 0.5,0.75,1.0,1.5,2.0\}\tabularnewline 
$\eta'$ & $5$ & \{ 1,3,5,10,20\} \tabularnewline
$Max_{it}$ & $\theta \times (|N'| + \sigma \times |T| + \phi \times |K|)$ & Maximum permissible number of consecutive iterations without improvement. \tabularnewline
$\theta$ & $3$ & Is redundant should be removed. \tabularnewline
$\sigma$ & $10$ &\{ 5,7,10,15,20\}. \tabularnewline
$\phi$ & $8$ & \{ 5,6,7,8,9,10,15,20\}. \tabularnewline
Disturb strength & $5$ & \{1,2,4,5,7,9,10\}\tabularnewline
\hline
\end{tabular}
\end{table}


% $Max_{it}$ & $\theta \times (|N'| + \sigma \times |T| + \phi \times |K|)$

$\theta$ = $3$ 
$\sigma$= $10$
$\phi$ = $10$ 
\begin{table}[H]
\caption{Summary of \hyperref[alg:MSILS]{MSILS} parameters}
\label{default parameters}
\centering
\begin{tabular}{ p{2cm}p{9cm}}
\toprule%
\textbf{Parameters}  &\textbf{Definition} \tabularnewline
\hline
$d_{\min}$& Minimum grid resolution in initial solution construction. \tabularnewline
$d_{\max}$& Increment defining the maximum possible grid resolution in initial solution construction. \tabularnewline
$Max_{restart}$  & Number of initial solutions. \tabularnewline
$U$  & Number that triggers the adjustment of the scalar $\eta$.\tabularnewline
$\kappa$  & The parameter for updaing the scalar $\eta$.\tabularnewline 
$\eta'$  & Initial value of $\eta$. \tabularnewline
$Max_{it}$  & Maximum permissible number of consecutive iterations without improvement. \tabularnewline
$\theta, \sigma, \phi$  & Parameters to compute $Max_{it}$. \tabularnewline
\hline
\end{tabular}
\end{table}
The proposed \hyperref[alg:MSILS]{MSILS} has several key components, including an application of the Lagrangian relaxation technique, the rebalance procedure, and the disturbance procedure. This section examines the performance gain from utilising each of these components. Four variants have been derived from the \hyperref[alg:MSILS]{MSILS}. Each variant deactivates or modifies a specific component. Table \ref{Variants of MSILS} summarises each variant.

\begin{table}\small
\caption{Variants derived from \hyperref[alg:MSILS]{MSILS}}
\label{Variants of MSILS}
\centering
\begin{tabular}{ p{1.5cm}p{10cm}}
\toprule%
\textbf{Variants} & \textbf{Description} \tabularnewline
\hline

\hyperref[alg:MSILS]{MSILS}& Full algorithm described in Section \ref{sec: alg}.  \tabularnewline
$\cancel{L}$ & Instead of initialise $\alpha$ as line \ref{LS:initialise alpha} in \hyperref[alg:MSILS]{MSILS}, $\alpha$ is set to 100000.\tabularnewline
$\cancel{R}$ & Deactivating the \hyperref[alg:rebalance]{Rebalance} Procedure at line \ref{run rebalance} in \hyperref[alg:MSILS]{MSILS}. \tabularnewline
$D\&R$ & Instead of using the \hyperref[alg:rebalance]{Rebalance} Procedure described in Section \ref{sec:rebalance}, line \ref{run rebalance} uses the \hyperref[alg: destroy and rebalance]{Destory\_And\_Rebalance} procedure described in Appendix~\ref{sec: rebalance without feasibility}\tabularnewline
$\cancel{D}$ & Deactivating the First-stage DISTURB from the procedure Perturbation. \tabularnewline
\hline
\end{tabular}
\end{table}

The variant $\cancel{L}$ is used to examine the performance gain of using the Lagrangian relaxation technique. By initialising $\alpha$ to a large value, the \hyperref[alg:LocalSearch]{Local\_Search} procedure would just repair the infeasible starting solution. Thereafter, due to the high penalty for stockout, the search process would remain in the feasible solution space. All four variants and the \hyperref[alg:MSILS]{MSILS} were tested on Archettie2012 benchmark. Each version of the algorithm was given a time limit of 1 minute. 


\subsection{Performance evaluation}\label{sec:performance evaluation}


Compared with the supplementary results from \cite{alvarez2018iterated} available at \url{https://www.dep.ufscar.br/munari/irp/}. 


The proposed \hyperref[alg:MSILS]{MSILS} is compared with ILS and SA in \cite{alvarez2018iterated}. The algorithm is validated in small instances by comparing the optimal solutions obtained by the ACS algorithm developed in \cite{archetti2019exact}.
\begin{table}[htbp] 
\centering 
\caption{Average percentage gap (\%) of \hyperref[alg:MSILS]{MSILS} compared with ILS} \label{tab:msils_vs_ils} 
\begin{tabular}{lccccc} 
\toprule & $|T|=2$ & $|T|=3$ & $|T|=4$ & $|T|=5$ & Average \\ 
\midrule 
$|N'|=50$ & -8.00 & -4.30 & -5.03 & -5.13 & -5.61 \\ 
$|N'|=100$ & -8.63 & -6.59 & -6.57 & -6.90 & -7.17 \\ 
$|N'|=200$ & -9.02 & -6.16 & -5.16 & -4.10 & -6.11 \\ 
\midrule 
Average & -8.55 & -5.68 & -5.58 & -5.38 & -6.30 \\ 
\bottomrule 
\end{tabular} 
\end{table}


\begin{table}[htbp] 
\centering 
\caption{Average percentage gap (\%) of \hyperref[alg:MSILS]{MSILS} compared with SA} \label{tab:msils_vs_ils} 
\begin{tabular}{lccccc} 
\toprule & $|T|=2$ & $|T|=3$ & $|T|=4$ & $|T|=5$ & Average \\ 
\midrule 
$|N'|=50$ & -7.59 & -5.28 & -4.82 & -4.97 & -5.67 \\ 
$|N'|=100$ & -10.00 & -6.66 & -7.56 & -6.16 & -7.59 \\ 
$|N'|=200$ & -15.59 & -10.18 & -9.77 & -7.40 & -10.73 \\ 
\midrule 
Average & -11.06 & -7.37 & -7.38 & -6.18 & -8.00 \\ 
\bottomrule 
\end{tabular} 
\end{table}


\subsection{\hyperref[alg:MSILS]{MSILS} with HGS} \label{comparions with HGS}

\appendix

\section{Functions used to construct initial solution}

\begin{algorithm}[H]
\captionsetup{labelformat=empty}
\caption{\textbf{Find Feasible Visit Time}($t_{current}$, $t_{Stockout}$, $i$)}
\label{alg:FindVisitTime}
\begin{algorithmic}[1]
\WHILE{true}
    \IF{$t_{current}\leq t_{Stockout}$}
        \STATE \textbf{// schedule forward//}
        \STATE $T_{candidate}\gets\{t\in T:  t_{current}\leq t \leq t_{stockout},$
        \STATE \hspace{9em}  customer $i$ is not visited at $t$, 
        \STATE \hspace{9em}  $Q\times |K| - \sum_{k\in K}\sum_{i\in N'} q^k_{i,t}(s) > 0\}$
        \IF{$T_{candidate}\neq [\hspace{0.1cm}]$}
        \RETURN $t \sim \textsc{Random}(T_{candidate})$
        \ELSE
            \STATE $t_{current} \gets t_{Stockout} +1$
        \ENDIF
    \ELSE
        \STATE \textbf{// schedule backward//}         
            \FOR{$t\gets t_{stockout}; t\geq 1;t\gets t-1$}
                \IF{customer $i$ is not visited at time period $t$}
                    \IF{$Q\times |K| - \sum_{k\in K}\sum_{i\in N'} q^k_{i,t}(s) > 0$}
                        \RETURN $t$
                    \ENDIF
                \ENDIF
            \ENDFOR
    \ENDIF
\ENDWHILE
\end{algorithmic}
\end{algorithm}




\begin{algorithm}[H]
\captionsetup{labelformat=empty}
\caption{\textbf{Find Vehicle To Deliver Max}($t$, $i$, $s$)}
\label{alg:FindVehicleToDeliverMax}
\begin{algorithmic}[1]
\STATE $k'\gets |K|+1$, $q'\gets 0$, 
\FOR{$k\in K$}
    \STATE $q\gets 0$
    \IF{$t=1$}
        \STATE $q\gets min\{Q-\sum_{j\in N'}q^k_{j,1}(s), U_{i} - I_{i,0}\}$    
    \ELSE
        \STATE $q\gets min\{Q-\sum_{j\in N'}q^k_{j,t}(s), U_{i} - I_{i,t-1}(s)\}$
    \ENDIF
    \IF{$q'<q$}
        \STATE $q' \gets q$, $k'\gets k$
    \ENDIF
\ENDFOR
\RETURN $k',q'$
\end{algorithmic}
\end{algorithm}




\section{\hyperref[alg:MSILS]{MSILS} with HGS}\label{MSILS with HGS}


\begin{algorithm}[H]
\caption{\textbf{\hyperref[alg:MSILS]{MSILS} with HGS in \cite{vidal2022hybrid}}}\label{alg:MSILS with HGS}
\begin{algorithmic}[1]
\STATE $s^* \gets$ Empty \AND  $f(s^*) \gets +\infty$  
\FOR{$i\gets 1$ \TO $Max_{restart}$} 
\STATE $s\gets$ \hyperref[alg:InitialSchedule]{Construct initial solution} 
\STATE $s \gets HGS(s)$ \COMMENT{to improve routing}
\IF{$f(s)<f(s^*)$} \STATE $s^*\gets s$ \ENDIF
\STATE Run\_HGS\_Or\_Not $\gets$ FALSE
\STATE  $it \gets 0$, $\eta \gets \eta'$, $u\gets 0$ 
\WHILE{ $it < Max_{it}$} 
    \STATE $\alpha \gets \epsilon (s^*, s , 0, \eta)$  
    \STATE $\bar{s} \gets s$
    \STATE $s\gets$ \hyperref[alg:LocalSearch]{Local\_Search}$(s,\alpha)$    
    \WHILE{$f_{\mu}(s,\alpha) \neq f_{\mu}(\bar{s},\alpha)$} 
        \STATE  $\alpha \gets \epsilon (s^*, s , \alpha, \eta)$       
        \STATE $\bar{s} \gets s$
        \STATE $s\gets$\hyperref[alg:LocalSearch]{Local\_Search}$(s,\alpha)$       
    \ENDWHILE
    \IF{$min\{it, u\} \geq U$}
        \STATE $u\gets 0$, $\eta \gets \kappa \times \eta$ 
    \ENDIF
    \IF{$s$ is feasible \AND $f(s)<f(s^*)$}
        \STATE $s^*\gets s$, $it\gets 0$
        \STATE $s^*\gets$ \hyperref[alg:rebalance]{Rebalance}$(s^*)$            
        \STATE Run\_HGS\_Or\_Not $\gets$ TRUE
    \ENDIF
    \STATE $s \gets$ Disturbance$(s^*, it)$, $it\gets it+1$, $u\gets u+1$ 
\ENDWHILE 
\IF{Run\_HGS\_Or\_Not = TRUE}
    \STATE $s^* \gets HGS(s^*)$
\ENDIF
\ENDFOR
\RETURN solution $s^*$
\end{algorithmic}
\end{algorithm}

\section{Functions used in Rebalance procedure guarantee feasible solution}



\begin{algorithm}[H]
\caption{\textbf{Feasible\_Quantity\_To\_Removal$(i,j)$}} \label{alg:FeasibleQuantityRemoval}
\begin{algorithmic}[1]
\STATE Quantity\_To\_Remove $\gets min\{I_{i,j}, q_{i,j}\}$ 
\STATE $I\gets I_{i,j-1}+ q_{i,j} -r_{i,j} - $Quantity\_To\_Remove \COMMENT{New inventory level of period $j$}
\STATE Stockout $\gets 0$
\FOR{$time\gets j $ \TO $|T|$}
\STATE $I\gets I_{i,time}+ q_{i,j} -r_{i,j}$
\IF{Stockout  $ \leq -I'$}
    \STATE Stockout $\gets -I'$
\ENDIF
\ENDFOR
\RETURN $\gets max\{0, Quantity - Stockout\}$
\end{algorithmic}
\end{algorithm}



\begin{algorithm}[H]
\caption{\textbf{Assign\_max\_possible\_quantity\_after\_feasible$(s)$}} \label{alg:AssignQuantity}
\begin{algorithmic}[1]

\FOR{$i \gets 1$ \TO $|T|$}     
\STATE $\text{PRIORITY} \gets \emptyset$
    
    \FOR{$j \gets 1$ \TO $|N'|$} 
        \IF{customer $j$ is visited on day $i$}
            \STATE Add customer $j$ to PRIORITY with computed priority value
        \ENDIF
    \ENDFOR
    
    \WHILE{PRIORITY is not empty}
        \STATE Select a customer randomly according to PRIORITY weights
        \STATE Deliver the maximum feasible quantity to the selected customer
        \STATE Remove the selected customer from PRIORITY
        \STATE Update inventory levels and delivery quantity for subsequent days
    \ENDWHILE
\ENDFOR
\RETURN $s$
\end{algorithmic}
\end{algorithm}


\section{Rebalance procedure without guarantee feasible solution}\label{sec: rebalance without feasibility}



\begin{algorithm}[H]
\caption{\textbf{Destory\_And\_Rebalance$(s^*)$}} \label{alg: destroy and rebalance}
\begin{algorithmic}[1]
\STATE $s \gets s^*$
\STATE Set all delivery quantities in $s$ to zero
\STATE Update inventory levels in $s$

\FOR{$i \gets 1$ \TO $|T|$}     
    \FOR{$j \gets 1$ \TO $|N'|$} 
        \IF{customer $j$ is visited on day $i$}
            \STATE Quantity $\gets $ \hyperref[alg:FindMaxQuantityAvoidStockout]{Find\_max\_quantity\_avoid\_stockout}$(j, i)$
            \STATE Assign Quantity
            \STATE Update inventory levels for all subsequent days
        \ENDIF
    \ENDFOR
\ENDFOR
\STATE Stockout\_resolved $\gets$ FALSE
\STATE $\{s,Stockout\_resolved\}\gets $\hyperref[alg:QunatityRepair]{Repair}$(s,Stockout\_resolved)$
\IF{Stock\_resolved = FALSE}
    \STATE return $s$
\ENDIF
\STATE $s\gets$Assign\_max\_possible\_quantity\_after\_feasible$(s)$
\STATE Remove visits with zero delivery quantities
\RETURN $s$
\end{algorithmic}
\end{algorithm}


\begin{algorithm}[H]
\caption{\textbf{Find\_max\_quantity\_avoid\_stockout$(customer, time)$}} \label{alg:FindMaxQuantityAvoidStockout}
\begin{algorithmic}[1]
\STATE Quantity $\gets 0$
\STATE Quantity $\gets$ \hyperref[alg:FindMaxStockout]{Find\_max\_stockout}$(customer, time)$
\STATE Quantity  $\gets \min\{Quantity, U_{customer} - \max\{I_{customer,time-1},0\}\}$ \COMMENT{If stockout occurs in the previous period, assumes the inventory level is 0.}
\STATE Quantity $\gets \min \{Quantity, Q - Vehicle\_load\} $
\RETURN Quantity
\end{algorithmic}
\end{algorithm}

\begin{algorithm}[H]
\caption{\textbf{Find\_max\_stockout$(customer, time)$}} \label{alg:FindMaxStockout}
\begin{algorithmic}[1]
\STATE Largest\_stockout $\gets 0$
\FOR{$t\gets time$ \TO $|T|$ }
    \IF{Largest\_stokout $> I_{customer, t}$}
        \STATE Largest\_stockout $\gets I_{customer,t}$ 
    \ENDIF
\ENDFOR
\STATE Largest\_stockout $\gets -Largest\_stockout$
\RETURN  Largest\_stockout
\end{algorithmic}
\end{algorithm}

\begin{algorithm}[H]
\caption{\textbf{Repair$(s, Stockout\_resolved)$}} \label{alg:QunatityRepair}
\begin{algorithmic}[1]
\STATE Possible\_repair $\gets$ TRUE
\WHILE{Stockout\_resolved = FALSE \AND Possible\_repair = TRUE}
    \STATE Stockout\_resolved $\gets$ TRUE
    \STATE Possible\_repair $\gets$ FALSE
    \FOR{$i \gets 1$ \TO $|T|$} 
        \FOR{$j \gets 1$ \TO $|N'|$}
            \IF{customer $j$ has stockout at time $i$}
                \STATE Stockout\_resolved $\gets$ FALSE
                \STATE $Quantity_{stockout} \gets I_{j,i}$
                \FOR{$k\gets 1$ \TO time $i$}
                    \IF{Customer $j$ is visited in time $k$}
                        \FOR{customer $z\in N'$ are visited by the same vehicle}
                            \IF{Some quantity, $\Delta$ associated with customer $z$ can be removed}
                                \FOR{$x \gets 1$ \TO $|T|$}
                                    \IF{customer $z$ is visited on time $x$}
                                        \IF{$\Delta$ can be delivered on time $x$ to customer $z$ without stockout}
                                            \STATE Possible\_repair $\gets$ TRUE
                                            \STATE $Quantity_{\Delta}\gets min\{Quantity_{stockout}, \Delta\}$
                                            \STATE Remove $Quantity_{\Delta}$ from visiting customer $z$ on time $k$
                                            \STATE Deliver extra $Quantity_{\Delta}$ unit on day $x$
                                            \STATE Deliver $Quantity_{\Delta}$ unit for customer $j$ on time $k$
                                            \STATE Update delivery quantity and inventory level 
                                        \ENDIF
                                    \ENDIF
                                \ENDFOR
                            \ENDIF
                        \ENDFOR                                    
                    \ENDIF
                \ENDFOR
            \ENDIF
        \ENDFOR
    \ENDFOR
\ENDWHILE
\RETURN $\{s,Stockout\_resolved\}$
\end{algorithmic}
\end{algorithm}

\bibliographystyle{apalike} 
\bibliography{reference}
\end{document}